\SourcePage{200}{203}
\section{Potential energy treated as a perturbation}

The case in which the particle's entire potential energy in an external field
can be treated as the perturbation calls for separate consideration. The
unperturbed Schrödinger equation is then the equation for free motion,
\begin{equation}
  \Delta\psi^{(0)}+k^2\psi^{(0)}=0,
  \qquad
  k=\frac{\sqrt{2mE}}{\hbar}=\frac p\hbar
  \tag{45.1}
\end{equation}
whose solutions are plane waves. Since the energy spectrum of free motion is
continuous, this is a special instance of perturbation theory in a continuous
spectrum. Here it is more convenient to obtain the solution directly, without
using the general formulas.

The equation for the first-order correction $\psi^{(1)}$ to the wave function
is
\begin{equation}
  \Delta\psi^{(1)}+k^2\psi^{(1)}=\frac{2mU}{\hbar^2}\psi^{(0)}
  \tag{45.2}
\end{equation}
($U$ is the potential energy). As is known from electrodynamics, its solution
may be written in the form of ``retarded potentials,'' namely\footnote{This is
a particular integral of equation (45.2). Any solution of the equation with
no right-hand side (that is, of the unperturbed equation (45.1)) may still be
added to it.}
\begin{equation}
  \psi^{(1)}(x,y,z)=-\frac{m}{2\pi\hbar^2}
  \int\psi^{(0)}U(x',y',z')e^{ikr}\frac{dV'}r,
  \tag{45.3}
\end{equation}
\[
  dV'=dx'\,dy'\,dz',
  \qquad
  r^2=(x-x')^2+(y-y')^2+(z-z')^2.
\]

Let us determine the conditions under which the field $U$ may be regarded as
a perturbation. Perturbation theory is applicable when
$\psi^{(1)}\ll\psi^{(0)}$. Let $a$ give the order of magnitude of the size of
the region in which the field differs appreciably from zero. Suppose first
that the particle's energy is so low that $ak$ is at most of order unity. For
an order-of-magnitude estimate the factor $e^{ikr}$ in the integrand of
(45.3) is then immaterial, and the whole integral is of order
$\psi^{(0)}|U|a^2$. Thus
$\psi^{(1)}\sim(ma^2|U|/\hbar^2)\psi^{(0)}$, and the condition becomes
\begin{equation}
  |U|\ll\frac{\hbar^2}{ma^2}
  \quad\text{when}\quad ka\lesssim1.
  \tag{45.4}
\end{equation}
\SourcePage{201}{204}
Notice that $\hbar^2/ma^2$ has a simple physical meaning: it gives the order
of magnitude of the kinetic energy that a particle confined to a volume of
linear dimensions $a$ would have (for its momentum would be
$\sim\hbar/a$ by the uncertainty relation).

As a particular example, consider a potential well so shallow that (45.4) is
satisfied. Such a well has no negative energy levels ($R.\,Peierls$, 1929),
as was already seen in the problem to \secref{33} for the special case of a
spherically symmetric well. Indeed, at $E=0$ the unperturbed wave function
reduces to a constant, which may conventionally be set equal to unity:
$\psi^{(0)}=1$. Since $\psi^{(1)}\ll\psi^{(0)}$, the wave function for motion
in the well, $\psi=1+\psi^{(1)}$, plainly vanishes nowhere. An eigenfunction
without nodes belongs to the normal state, so $E=0$ remains the lowest
possible energy of the particle.

Thus, if a well is not sufficiently deep, only infinite motion of the
particle is possible: the well cannot ``capture'' it. This result is
specifically quantum-mechanical. In classical mechanics a particle can
undergo finite motion in any potential well.

It must be stressed that everything just said applies only to a
three-dimensional well. A one- or two-dimensional well (that is, one whose
field depends on only one or two coordinates) always has negative energy
levels (see the problems to this section). The reason is that in one and two
dimensions the perturbation theory under consideration is wholly inapplicable
at zero (or very small) energy $E$.
\footnote{In two dimensions, $\psi^{(1)}$ is expressed (as follows from the
theory of the two-dimensional wave equation) by an integral analogous to
(45.3), with $i\pi H_0^{(1)}(kr)dx'\,dy'$ in place of
$e^{ikr}dx'\,dy'\,dz'/r$ ($H_0^{(1)}$ is a Hankel function), and
$r=\sqrt{(x'-x)^2+(y'-y)^2}$. As $k\to0$, the Hankel function, and with it
the entire integral, tends logarithmically to infinity.

Similarly, in one dimension the integrand defining $\psi^{(1)}$ contains
$2\pi i e^{ikr}dx'/k$ (where $r=|x'-x|$), and as $k\to0$,
$\psi^{(1)}$ tends to infinity as $1/k$.}

At high energies, where $ka\gg1$, the factor $e^{ikr}$ in the integrand is
essential and greatly reduces the magnitude of the integral. In this case
solution (45.3) can be
\SourcePage{202}{205}
put into another form. For its derivation, however, it is simpler to return
directly to equation (45.2). Choose the direction of unperturbed motion as the
$x$ axis. The unperturbed wave function is then $\psi^{(0)}=e^{ikx}$ (the
constant factor is conventionally set to unity). Seek a solution of
\[
  \Delta\psi^{(1)}+k^2\psi^{(1)}=\frac{2m}{\hbar^2}Ue^{ikx}
\]
in the form $\psi^{(1)}=e^{ikx}f$. Since $k$ is assumed large, in
$\Delta\psi^{(1)}$ it suffices to retain only terms in which the factor
$e^{ikx}$ is differentiated at least once. The equation for $f$ is then
\[
  2ik\frac{\partial f}{\partial x}=\frac{2mU}{\hbar^2},
\]
and hence
\begin{equation}
  \psi^{(1)}=e^{ikx}f=-\frac{im}{\hbar^2k}e^{ikx}\int U\,dx.
  \tag{45.5}
\end{equation}

The integral is of order
$|\psi^{(1)}|\sim m|U|a/\hbar^2k$, so in this case perturbation theory is
applicable under the condition
\begin{equation}
  |U|\ll\frac{\hbar^2}{ma^2}ka=\frac{\hbar v}{a},
  \qquad ka\gg1
  \tag{45.6}
\end{equation}
($v=k\hbar/m$ is the particle's velocity). This condition is weaker than
(45.4). Consequently, if the field may be treated as a perturbation at low
particle energies, it certainly may be so treated at high energies; the
converse, generally speaking, does not hold.
\footnote{In one dimension the condition for applicability of perturbation
theory is the inequality (45.6) for all $ka$. The derivation of (45.4) given
above for three dimensions cannot be carried out in one dimension because of
the divergence, noted in the footnote on p.~204, of the function
$\psi^{(1)}$ constructed in that manner.}

The applicability of the perturbation theory developed here to the Coulomb
field requires separate consideration. In the field $U=\alpha/r$ there is no
finite region of space outside which $U$ is substantially smaller than inside
it. The required condition follows by replacing the parameter $a$ in (45.6)
with the variable distance $r$; this gives
\begin{equation}
  \frac\alpha{\hbar v}\ll1.
  \tag{45.7}
\end{equation}
\SourcePage{203}{206}
Thus the Coulomb field may be treated as a perturbation at high particle
energies.
\footnote{It should be borne in mind that the integral (45.5) with
$U=\alpha/r$ diverges logarithmically at large
$x/\sqrt{y^2+z^2}$. The Coulomb-field wave function obtained by perturbation
theory is therefore inapplicable inside a narrow cone about the $x$ axis.}

Finally, let us derive a formula giving an approximate wave function for a
particle whose energy $E$ is everywhere much greater than the potential
energy $U$ (no other conditions are required). To the first approximation,
the coordinate dependence of the wave function is the same as in free motion,
whose direction we choose as the $x$ axis. Accordingly, write
$\psi=e^{ikx}F$, where $F$ is a function of the coordinates that varies slowly
in comparison with the factor $e^{ikx}$ (although in general it cannot be
said to be close to unity). Substitution in the Schrödinger equation gives
\begin{equation}
  2ik\frac{\partial F}{\partial x}=\frac{2m}{\hbar^2}UF,
  \tag{45.8}
\end{equation}
whence
\begin{equation}
  \psi=e^{ikx}F=\operatorname{const}\cdot e^{ikx}
  \exp\left(-\frac{i}{\hbar v}\int U\,dx\right).
  \tag{45.9}
\end{equation}
This is the required expression. It must be remembered, however, that it is
not valid at excessively large distances. In (45.8) the term $\Delta F$,
which contains second derivatives of $F$, was omitted. At large distances
$\partial^2F/\partial x^2$ tends to zero together with
$\partial F/\partial x$. The derivatives with respect to the transverse
coordinates $y,z$, on the other hand, do not tend to zero, and they may be
neglected only if $x\ll ka^2$.

\ProblemHeading
\textbf{1.} Determine the energy level in a shallow one-dimensional
potential well; condition (45.4) is assumed to hold.

\SolutionHeading Make the assumption, confirmed by the result,
that $|E|\ll|U|$. Then $E$ may be neglected on the right-hand side of the
Schrödinger equation within the well,
\[
  \frac{d^2\psi}{dx^2}=\frac{2m}{\hbar^2}(U(x)-E)\psi,
\]
and $\psi$ may also be treated as a constant, which can be set equal to unity
without loss of generality:
\[
  \frac{d^2\psi}{dx^2}=\frac{2m}{\hbar^2}U.
\]
\SourcePage{204}{207}
Integrate this equality with respect to $x$ between two points $\pm x_1$ such
that $a\ll x_1\ll1/\kappa$, where $a$ is the width of the well and
$\kappa=\sqrt{2m|E|}/\hbar$. Since the integral of $U(x)$ converges, the
integration on the right may be extended over the whole interval from
$-\infty$ to $+\infty$:
\begin{equation}
  \left.\frac{d\psi}{dx}\right|_{-x_1}^{x_1}
  =\frac{2m}{\hbar^2}\int_{-\infty}^{+\infty}U\,dx.
  \tag{1}
\end{equation}
Far from the well the wave function has the form $\psi=e^{\mp\kappa x}$.
Substitution in (1) gives
\[
  -2\kappa=\frac{2m}{\hbar^2}\int_{-\infty}^{+\infty}U\,dx
  \qquad\text{or}\qquad
  |E|=\frac{m}{2\hbar^2}
  \left(\int_{-\infty}^{+\infty}U\,dx\right)^2.
\]
In accord with the assumption made, the magnitude of the level is thus a
small quantity of higher (second) order than the depth of the well.

\ProblemHeading
\textbf{2.} Determine the energy level in a shallow two-dimensional
potential well $U(r)$ ($r$ is the polar coordinate in the plane); the integral
$\int_0^\infty rU\,dr$ is assumed to converge.

\SolutionHeading Proceeding as in the preceding problem, within
the well we obtain
\[
  \frac1r\frac d{dr}\left(r\frac{d\psi}{dr}\right)
  =\frac{2m}{\hbar^2}U.
\]
Integrating with respect to $r$ from 0 to $r_1$ (where
$a\ll r_1\ll1/\kappa$), we find
\begin{equation}
  \left.\frac{d\psi}{dr}\right|_{r=r_1}
  =\frac{2m}{\hbar^2r_1}\int_0^\infty rU(r)\,dr.
  \tag{1}
\end{equation}
Far from the well, the two-dimensional free-motion equation
\[
  \frac1r\frac d{dr}\left(r\frac{d\psi}{dr}\right)
  +\frac{2m}{\hbar^2}E\psi=0
\]
has the solution (vanishing at infinity)
$\psi=\operatorname{const}\cdot H_0^{(1)}(i\kappa r)$. At small values of
the argument, the leading term of this function is proportional to
$\ln\kappa r$. Bearing this in mind, at $r\sim a$ equate the logarithmic
derivatives of $\psi$ evaluated inside the well (the right-hand side of (1))
and outside it. This gives
\[
  -\frac1{a\ln\kappa a}\approx
  \frac{2m}{\hbar^2a}\int_0^\infty U(r)r\,dr,
\]
and hence
\[
  |E|\sim\frac{\hbar^2}{ma^2}
  \exp\left[-\frac{\hbar^2}{m}
  \left|\int_0^\infty Ur\,dr\right|^{-1}\right].
\]
Thus the energy level is exponentially small in comparison with the depth of
the well.
